\chapter[Metodi Diretti per Sistemi Lineari e Fattorizzazione LU]{Metodi Diretti per Sistemi Lineari e Fattorizzazione LU:\\Algoritmi, Stabilit\`a Numerica e Pivoting}
\label{chap:metodi-diretti-fattorizzazione-lu}

\FloatBarrier
\section{Definizione di Sistema Lineare e Geometria delle Soluzioni}
\label{sec:def-sistema-lineare-geometria}

\subsection{Formulazione Algebrica Compatta}
Un sistema di $m$ equazioni lineari in $n$ incognite a coefficienti reali si esprime in forma matriciale compatta come:
\begin{equation}
    Ax = b,
    \label{eq:sistema-lineare-generale}
\end{equation}
in cui:
\begin{itemize}
    \item $A \in \R^{m \times n}$ \`e la matrice dei coefficienti del sistema, nota;
    \item $b \in \R^m$ \`e il vettore dei termini noti, noto;
    \item $x \in \R^n$ \`e il vettore delle incognite da determinare.
\end{itemize}

Risolvere il sistema~\eqref{eq:sistema-lineare-generale} equivale a caratterizzare la totalit\`a dei vettori $x \in \R^n$ che rendono vera l'uguaglianza vettoriale.

\subsection{Interpretazione Algebrico-Geometrica: Combinazione Lineare di Colonne}
Esprimendo la matrice $A$ per colonne, ovvero denotando con $a_j \in \R^m$ la $j$-esima colonna di $A$ per $j = 1, \dots, n$:
\begin{equation}
    A = \begin{bmatrix}
        \vert & \vert & & \vert \\
        a_1 & a_2 & \dots & a_n \\
        \vert & \vert & & \vert
    \end{bmatrix}, \qquad
    x = \begin{bmatrix}
        x_1 \\ x_2 \\ \vdots \\ x_n
    \end{bmatrix},
\end{equation}
l'operazione di moltiplicazione matrice-vettore si traduce identicamente nella combinazione lineare:
\begin{equation}
    Ax = \sum_{j=1}^n x_j a_j = x_1 a_1 + x_2 a_2 + \dots + x_n a_n.
    \label{eq:combinazione-lineare-colonne}
\end{equation}

\begin{property}[Interpretazione Geometrica di Compatibilit\`a]
La risoluzione di $Ax = b$ equivale a determinare i coefficienti scalari $x_1, \dots, x_n \in \R$ tali che la combinazione lineare delle colonne di $A$ riproduca esattamente il vettore $b$. Ne conseguono le seguenti tre situazioni mutualmente esclusive:
\begin{enumerate}
    \item \textbf{Sistema Incompatibile (Nessuna soluzione):} Il vettore $b$ non appartiene allo spazio generato dalle colonne di $A$:
    \begin{equation}
        b \notin \mathrm{Range}(A) \equiv \spanop(a_1, a_2, \dots, a_n).
    \end{equation}
    \item \textbf{Sistema Compatibile Indeterminato (Infinite soluzioni):} Il vettore appartiene all'immagine ($b \in \mathrm{Range}(A)$), ma le colonne di $A$ sono linearmente dipendenti, ovvero il nucleo di $A$ non \`e degenere:
    \begin{equation}
        \dim(\ker(A)) = n - \rank(A) > 0.
    \end{equation}
    \item \textbf{Sistema Compatibile Determinato (Soluzione unica):} Il vettore appartiene all'immagine ($b \in \mathrm{Range}(A)$) e le colonne di $A$ sono linearmente indipendenti, ovvero $\ker(A) = \{0\}$.
\end{enumerate}
\end{property}

\subsection{Esempi Notabili: Struttura del Nucleo e Compatibilit\`a}

\begin{example}[Sistema Indeterminato con Infinite Soluzioni]
\label{ex:infinite-soluzioni}
Si consideri la matrice $A \in \R^{4 \times 3}$ e il vettore $b \in \R^4$ definiti da:
\begin{equation}
    A = \begin{bmatrix}
        2 & 0 & 1 \\
        0 & 1 & 1 \\
        0 & 0 & 0 \\
        0 & 0 & 0
    \end{bmatrix}, \qquad
    b = \begin{bmatrix}
        4 \\ 4 \\ 0 \\ 0
    \end{bmatrix}.
\end{equation}
Il vettore delle incognite $x$ deve necessariamente appartenere a $\R^3$ ($n = 3$, pari al numero di colonne di $A$). Esaminando il sistema per componenti:
\begin{equation}
    \begin{cases}
        2x_1 + x_3 = 4 \\
        x_2 + x_3 = 4 \\
        0 = 0 \\
        0 = 0
    \end{cases}
\end{equation}
Si possono individuare immediatamente pi\`u vettori soluzione, ad esempio:
\begin{equation}
    x^{(1)} = \begin{bmatrix} 1 \\ 2 \\ 2 \end{bmatrix}, \qquad
    x^{(2)} = \begin{bmatrix} 0 \\ 0 \\ 4 \end{bmatrix}, \qquad
    x^{(3)} = \begin{bmatrix} 2 \\ 4 \\ 0 \end{bmatrix}.
\end{equation}
Verifichiamo analiticamente la prima soluzione $x^{(1)}$:
\begin{equation}
    A x^{(1)} = \begin{bmatrix}
        2(1) + 0(2) + 1(2) \\
        0(1) + 1(2) + 1(2) \\
        0 \\ 0
    \end{bmatrix} = \begin{bmatrix} 4 \\ 4 \\ 0 \\ 0 \end{bmatrix} = b.
\end{equation}
Per dimostrare l'esistenza di infinite soluzioni, calcoliamo il nucleo $\ker(A) = \{w \in \R^3 \mid Aw = 0\}$:
\begin{equation}
    \begin{cases}
        2w_1 + w_3 = 0 \implies w_1 = -\frac{1}{2}w_3 \\
        w_2 + w_3 = 0 \implies w_2 = -w_3
    \end{cases}
    \implies
    w = \alpha \begin{bmatrix} -1/2 \\ -1 \\ 1 \end{bmatrix} \quad (\alpha \in \R).
\end{equation}
Dato un qualsiasi vettore $w \in \ker(A)$ e una soluzione particolare $x$, per linearit\`a il vettore $\tilde{x} = x + w$ soddisfa:
\begin{equation}
    A\tilde{x} = A(x + w) = Ax + Aw = b + 0 = b.
\end{equation}
Poich\'e $\dim(\ker(A)) = 3 - 2 = 1$, la variet\`a lineare delle soluzioni \`e una retta in $\R^3$:
\begin{equation}
    \mathcal{S} = \left\{ \begin{bmatrix} 2 \\ 4 \\ 0 \end{bmatrix} + \alpha \begin{bmatrix} -1 \\ -2 \\ 2 \end{bmatrix} \;\middle|\; \alpha \in \R \right\}.
\end{equation}
\end{example}

\begin{example}[Sistema Incompatibile]
\label{ex:sistema-impossibile}
Mantenendo inalterata la matrice $A$ dell'Esempio~\ref{ex:infinite-soluzioni}, si scelga il termine noto:
\begin{equation}
    b = \begin{bmatrix} 4 \\ 4 \\ 1 \\ 0 \end{bmatrix}.
\end{equation}
Poich\'e le ultime due righe di $A$ sono interamente nulle, per ogni possibile vettore $x \in \R^3$ la terza e la quarta componente del prodotto $Ax$ valgono identicamente zero:
\begin{equation}
    Ax = \begin{bmatrix} * \\ * \\ 0 \\ 0 \end{bmatrix} \neq \begin{bmatrix} 4 \\ 4 \\ 1 \\ 0 \end{bmatrix} = b.
\end{equation}
La terza equazione si ridurrebbe a $0 = 1$, assurdo. Geometricamente, $b \notin \mathrm{Range}(A)$, rendendo il sistema incompatibile.
\end{example}

\FloatBarrier
\section{Sistemi Quadrati, Invertibilit\`a e la ``Regola Aurea''}
\label{sec:matrici-invertibili-regola-aurea}

\subsection{Matrici Non Singolari e Unicit\`a Globale}
Quando il sistema presenta lo stesso numero di equazioni e di incognite ($m = n$), la matrice $A \in \R^{n \times n}$ \`e quadrata.

\begin{theorem}[Esistenza e Unicit\`a per Matrici Invertibili]
\label{thm:esistenza-unicita-quadrata}
Sia $A \in \R^{n \times n}$ una matrice quadrata. Le seguenti condizioni sono logicamente equivalenti:
\begin{enumerate}
    \item $A$ \`e invertibile (o non singolare), ovvero esiste $A^{-1} \in \R^{n \times n}$ tale che $A A^{-1} = A^{-1} A = I_n$;
    \item $\det(A) \neq 0$;
    \item $\rank(A) = n$;
    \item $\ker(A) = \{0\}$;
    \item Le $n$ colonne di $A$ (cos\`i come le $n$ righe) formano una base di $\R^n$.
\end{enumerate}
Se tali condizioni sono soddisfatte, allora per \textbf{qualsiasi} vettore termine noto $b \in \R^n$, il sistema $Ax = b$ ammette una e una sola soluzione:
\begin{equation}
    x = A^{-1}b.
    \label{eq:soluzione-formale-inversa}
\end{equation}
\end{theorem}

\subsection{La Regola Aurea del Calcolo Numerico: Mai Calcolare $A^{-1}$}
\label{subsec:regola-aurea}

Sebbene l'equazione~\eqref{eq:soluzione-formale-inversa} sia matematicamente ineccepibile, \textbf{tradurla pedissequamente in codice numerico calcolando esplicitamente $A^{-1}$ ed effettuando poi il prodotto $x = A^{-1}b$ costituisce un grave errore computazionale}~\cite{higham2002accuracy,trefethen1997numerical}.

\begin{alertblock}{Regola Aurea dell'Algebra Lineare Numerica}
Non calcolare mai esplicitamente l'inversa $A^{-1}$ di una matrice per risolvere un sistema lineare $Ax = b$.
\end{alertblock}

Le motivazioni sono duplici ed entrambe decisive:
\begin{enumerate}
    \item \textbf{Costo Computazionale Superfluo:} Calcolare esplicitamente $A^{-1}$ richiede l'equivalente della risoluzione di $n$ sistemi lineari distinti (aventi come termini noti le $n$ colonne della matrice identit\`a $I_n$), comportando circa $2n^3$ operazioni floating-point (FLOPs). Al contrario, fattorizzare la matrice e risolvere per sostituzione richiede soltanto $\frac{2}{3}n^3$ FLOPs: calcolare l'inversa \`e esattamente tre volte pi\`u lento.
    \item \textbf{Instabilit\`a Numerica e Perdita di Precisione:} L'inversione esplicita introduce errori di arrotondamento significativamente superiori rispetto ai metodi basati su fattorizzazione diretta ($LU$ o $QR$). Nel modello a virgola mobile, la matrice computata $\mathrm{fl}(A^{-1})$ pu\`o distorcere sensibilmente la direzione della soluzione se il numero di condizionamento $\kappa(A) = \|A\| \|A^{-1}\|$ \`e elevato.
    \item \textbf{Distruzione della Sparsit\`a:} Se $A$ \`e una matrice sparsa di grandi dimensioni (ad esempio tridiagonale o a bande, con $\mathcal{O}(n)$ elementi non nulli), la sua inversa $A^{-1}$ \`e quasi invariabilmente \textbf{densa}, richiedendo $\mathcal{O}(n^2)$ memoria e rendendo impossibile la scalabilit\`a su calcolatori moderni.
\end{enumerate}

Nella prassi computazionale, si ricorre sempre a solutori integrati che sfruttano fattorizzazioni interne:
\begin{lstlisting}[style=pythoncode, caption={Risoluzione corretta in Python mediante SciPy e confronto concettuale.}, label={lst:scipy-solve}]
import numpy as np
import scipy.linalg

# Generazione di matrice e termine noto
n = 1000
A = np.random.randn(n, n)
b = np.random.randn(n)

# APPROCCIO CORRETTO (fattorizzazione LU interna in O(2/3 n^3))
x_corretto = scipy.linalg.solve(A, b)

# APPROCCIO DA EVITARE ASSOLUTAMENTE (spreco O(2 n^3) e instabile)
# x_errato = np.linalg.inv(A) @ b
\end{lstlisting}

In ambiente MATLAB e Octave, l'operatore deputato \`e il backslash (\texttt{mldivide}):
\begin{lstlisting}[style=code, caption={Sintassi standard in MATLAB/Octave.}, label={lst:matlab-backslash}]
% Operatore mldivide: determina automaticamente la fattorizzazione ottimale
x = A \ b;
\end{lstlisting}

\FloatBarrier
\section{I Quattro Casi Facili (Easy Cases)}
\label{sec:easy-cases}

Prima di affrontare matrici dense generali mediante eliminazione gaussiana, l'algebra lineare numerica analizza quattro classi di matrici strutturate per le quali il sistema $Ax = b$ ammette una risoluzione immediata ed economica.

\subsection{Caso 1: Matrice Diagonale}
Sia $A = \diag(a_{11}, a_{22}, \dots, a_{nn})$ con elementi diagonali tutti non nulli ($a_{ii} \neq 0$ per $i = 1, \dots, n$):
\begin{equation}
    \begin{bmatrix}
        a_{11} & 0 & \dots & 0 \\
        0 & a_{22} & \dots & 0 \\
        \vdots & \vdots & \ddots & \vdots \\
        0 & 0 & \dots & a_{nn}
    \end{bmatrix}
    \begin{bmatrix}
        x_1 \\ x_2 \\ \vdots \\ x_n
    \end{bmatrix}
    =
    \begin{bmatrix}
        b_1 \\ b_2 \\ \vdots \\ b_n
    \end{bmatrix}.
\end{equation}
Il sistema si disaccoppia in $n$ equazioni scalari indipendenti:
\begin{equation}
    a_{ii} x_i = b_i \iff x_i = \frac{b_i}{a_{ii}}, \qquad \forall i = 1, \dots, n.
    \label{eq:soluzione-diagonale}
\end{equation}
\begin{property}[Costo e Memoria per Matrici Diagonali]
La risoluzione richiede esattamente $n$ divisioni scalari, con costo computazionale $\mathcal{O}(n)$. In memoria, non si memorizza mai la matrice $n \times n$ (che sprecherebbe $\mathcal{O}(n^2)$ spazio per contenere quasi esclusivamente zeri), bens\`i il solo vettore diagonale $d \in \R^n$.
\end{property}

\subsection{Caso 2: Matrice Triangolare Inferiore (Forward Substitution)}
Sia $L \in \R^{n \times n}$ una matrice triangolare inferiore non singolare ($l_{ii} \neq 0$ per ogni $i$):
\begin{equation}
    L = \begin{bmatrix}
        l_{11} & 0 & 0 & \dots & 0 \\
        l_{21} & l_{22} & 0 & \dots & 0 \\
        l_{31} & l_{32} & l_{33} & \dots & 0 \\
        \vdots & \vdots & \vdots & \ddots & \vdots \\
        l_{n1} & l_{n2} & l_{n3} & \dots & l_{nn}
    \end{bmatrix}.
\end{equation}
Scrivendo il sistema $Lx = b$ per esteso:
\begin{equation}
    \begin{aligned}
        l_{11} x_1 &= b_1 \\
        l_{21} x_1 + l_{22} x_2 &= b_2 \\
        l_{31} x_1 + l_{32} x_2 + l_{33} x_3 &= b_3 \\
        &\;\;\vdots \\
        l_{n1} x_1 + l_{n2} x_2 + \dots + l_{nn} x_n &= b_n.
    \end{aligned}
\end{equation}

L'algoritmo di \textbf{sostituzione in avanti} (\textit{Forward Substitution}) procede sequenzialmente dall'alto verso il basso:
\begin{enumerate}
    \item Dalla riga 1:
    \begin{equation}
        x_1 = \frac{b_1}{l_{11}}.
    \end{equation}
    \item Sostituendo $x_1$ nella riga 2, si isola l'incognita $x_2$:
    \begin{equation}
        x_2 = \frac{b_2 - l_{21}x_1}{l_{22}}.
    \end{equation}
    \item Proseguendo induttivamente, per la riga generica $i \in \{1, \dots, n\}$:
    \begin{equation}
        x_i = \frac{b_i - \sum_{j=1}^{i-1} l_{ij} x_j}{l_{ii}}.
        \label{eq:formula-forward-sub}
    \end{equation}
\end{enumerate}

\begin{algorithm}[H]
\caption{Forward Substitution (Sostituzione in Avanti)}
\label{alg:forward-substitution}
\KwInput{Matrice triangolare inferiore $L \in \R^{n \times n}$ con $l_{ii} \neq 0$, vettore termine noto $b \in \R^n$}
\KwOutput{Vettore soluzione $x \in \R^n$}
\For{$i = 1$ \KwTo $n$}{
    $s \leftarrow 0$\;
    \For{$j = 1$ \KwTo $i - 1$}{
        $s \leftarrow s + l_{ij} \cdot x_j$\;
    }
    $x_i \leftarrow \frac{b_i - s}{l_{ii}}$\;
}
\Return $x$\;
\end{algorithm}

\subsubsection{Conteggio delle Operazioni e Dimostrazione Geometrica}
Contiamo il numero esatto di operazioni floating-point (moltiplicazioni, sottrazioni e divisioni):
\begin{itemize}
    \item Riga $1$: $1$ divisione $\implies 1$ operazione;
    \item Riga $2$: $1$ moltiplicazione, $1$ sottrazione, $1$ divisione $\implies 3$ operazioni;
    \item Riga $3$: $2$ moltiplicazioni, $2$ sottrazioni, $1$ divisione $\implies 5$ operazioni;
    \item Riga $i$: $i - 1$ prodotti, $i - 1$ sottrazioni, $1$ divisione $\implies 2(i - 1) + 1 = 2i - 1$ operazioni;
    \item Riga $n$: $2n - 1$ operazioni.
\end{itemize}

Il costo totale \`e dato dalla somma dei primi $n$ numeri dispari:
\begin{equation}
    \text{FLOPs} = \sum_{i=1}^n (2i - 1) = 1 + 3 + 5 + \dots + (2n - 1).
\end{equation}

\begin{figure}[htbp]
\centering
\begin{tikzpicture}[scale=0.72, every node/.style={transform shape}]
    % Griglia n x n
    \draw[step=1cm, gray!25, very thin] (0,0) grid (5,5);
    \draw[thick, black!80] (0,0) rectangle (5,5);

    % Gnomone 1 (blu scuro)
    \fill[blue!70] (0,4) rectangle (1,5);
    \node at (0.5, 4.5) {\textbf{\footnotesize 1}};

    % Gnomone 2 (azzurro) - 3 quadrati
    \fill[cyan!60] (1,4) rectangle (2,5);
    \fill[cyan!60] (0,3) rectangle (2,4);
    \node at (1.5, 4.5) {\footnotesize 3};
    \node at (0.5, 3.5) {\footnotesize 3};
    \node at (1.5, 3.5) {\footnotesize 3};

    % Gnomone 3 (verde) - 5 quadrati
    \fill[teal!50] (2,3) rectangle (3,5);
    \fill[teal!50] (0,2) rectangle (3,3);
    \node at (2.5, 4.0) {\footnotesize 5};
    \node at (1.0, 2.5) {\footnotesize 5};

    % Gnomone 4 (arancione) - 7 quadrati
    \fill[orange!50] (3,2) rectangle (4,5);
    \fill[orange!50] (0,1) rectangle (4,2);
    \node at (3.5, 3.5) {\footnotesize 7};
    \node at (1.5, 1.5) {\footnotesize 7};

    % Gnomone 5 (rosso) - 9 quadrati
    \fill[red!40] (4,1) rectangle (5,5);
    \fill[red!40] (0,0) rectangle (5,1);
    \node at (4.5, 3.0) {\footnotesize 9};
    \node at (2.0, 0.5) {\footnotesize 9};

    % Indicazioni dimensionali
    \draw[<->, thick, navy] (0, 5.3) -- (5, 5.3) node[midway, above] {$n$};
    \draw[<->, thick, navy] (-0.3, 0) -- (-0.3, 5) node[midway, left] {$n$};
    
    \node[anchor=west] at (6.0, 3.2) {\textbf{Dimostrazione Visiva (Gnomoni Pitagorici):}};
    \node[anchor=west] at (6.0, 2.3) {Passo 1: $1 = 1^2$};
    \node[anchor=west] at (6.0, 1.6) {Passo 2: $1 + 3 = 4 = 2^2$};
    \node[anchor=west] at (6.0, 0.9) {Passo 3: $1 + 3 + 5 = 9 = 3^2$};
    \node[anchor=west] at (6.0, 0.2) {Passo $n$: $\sum_{i=1}^n (2i-1) = n^2$};
\end{tikzpicture}
\caption{Dimostrazione geometrica visiva della somma dei primi $n$ numeri dispari mediante accrescimento gnomonico: l'unione ordinata a forma di ``L'' produce esattamente un'area quadrata di lato $n$, pari a $n^2$.}
\label{fig:somma-dispari-gnomoni}
\end{figure}

Dunque il costo esatto \`e:
\begin{equation}
    \sum_{i=1}^n (2i - 1) = n^2 \implies \mathcal{O}(n^2) \text{ operazioni.}
\end{equation}
Poich\'e una matrice triangolare contiene $\frac{n(n+1)}{2} \approx \frac{n^2}{2}$ elementi non nulli che devono essere necessariamente processati, una complessit\`a pari a $\mathcal{O}(n^2)$ \`e computazionalmente \textbf{ottimale}.

\subsection{Caso 3: Matrice Triangolare Superiore (Backward Substitution)}
Sia $U \in \R^{n \times n}$ una matrice triangolare superiore non singolare ($u_{ii} \neq 0$ per ogni $i$):
\begin{equation}
    U = \begin{bmatrix}
        u_{11} & u_{12} & u_{13} & \dots & u_{1n} \\
        0 & u_{22} & u_{23} & \dots & u_{2n} \\
        0 & 0 & u_{33} & \dots & u_{3n} \\
        \vdots & \vdots & \vdots & \ddots & \vdots \\
        0 & 0 & 0 & \dots & u_{nn}
    \end{bmatrix}.
\end{equation}
L'algoritmo di \textbf{sostituzione all'indietro} (\textit{Backward Substitution}) procede a ritroso, partendo dall'ultima riga e risalendo fino alla prima:
\begin{enumerate}
    \item Dall'ultima equazione ($n$-esima):
    \begin{equation}
        u_{nn} x_n = b_n \implies x_n = \frac{b_n}{u_{nn}}.
    \end{equation}
    \item Dalla penultima equazione ($(n-1)$-esima), sostituendo $x_n$:
    \begin{equation}
        x_{n-1} = \frac{b_{n-1} - u_{n-1, n} x_n}{u_{n-1, n-1}}.
    \end{equation}
    \item Per la generica riga $i$, procedendo da $i = n-1$ decrescendo fino a $1$:
    \begin{equation}
        x_i = \frac{b_i - \sum_{j=i+1}^n u_{ij} x_j}{u_{ii}}.
        \label{eq:formula-backward-sub}
    \end{equation}
\end{enumerate}

\begin{algorithm}[H]
\caption{Backward Substitution (Sostituzione all'Indietro)}
\label{alg:backward-substitution}
\KwInput{Matrice triangolare superiore $U \in \R^{n \times n}$ con $u_{ii} \neq 0$, vettore termine noto $b \in \R^n$}
\KwOutput{Vettore soluzione $x \in \R^n$}
\For{$i = n, n-1, \dots, 1$}{
    $s \leftarrow 0$\;
    \For{$j = i + 1$ \KwTo $n$}{
        $s \leftarrow s + u_{ij} \cdot x_j$\;
    }
    $x_i \leftarrow \frac{b_i - s}{u_{ii}}$\;
}
\Return $x$\;
\end{algorithm}

Per simmetria rispetto al caso triangolare inferiore, il numero di operazioni \`e esattamente identico:
\begin{equation}
    \text{FLOPs} = \sum_{k=1}^n (2k - 1) = n^2 \implies \mathcal{O}(n^2) \text{ operazioni.}
\end{equation}

\subsection{Caso 4: Matrice Ortogonale}
Una matrice quadrata $Q \in \R^{n \times n}$ si dice \textbf{ortogonale} se le sue colonne formano una base ortonormale di $\R^n$, propriet\`a sintetizzata dalla relazione:
\begin{equation}
    Q^T Q = Q Q^T = I_n \iff Q^{-1} = Q^T.
\end{equation}
Per risolvere il sistema $Qx = b$, \`e sufficiente moltiplicare a sinistra per la trasposta $Q^T$:
\begin{equation}
    Q^T Q x = Q^T b \implies I_n x = Q^T b \implies x = Q^T b.
\end{equation}
Non \`e richiesta alcuna eliminazione o sostituzione: la soluzione coincide con il prodotto matrice-vettore tra la trasposta e il termine noto.
\begin{property}[Complessit\`a del Caso Ortogonale]
Il calcolo di $x = Q^T b$ equivale a calcolare gli $n$ prodotti scalari tra le colonne di $Q$ e il vettore $b$. Ciascun prodotto scalare richiede $n$ moltiplicazioni e $n-1$ somme: il costo computazionale totale \`e di $2n^2 - n \approx 2n^2$ FLOPs, ovvero $\mathcal{O}(n^2)$.
\end{property}

\FloatBarrier
\section{Metodi Diretti vs Metodi Iterativi}
\label{sec:diretti-vs-iterativi}

Quando la matrice dei coefficienti $A$ non appartiene a nessuna delle quattro classi strutturate analizzate, il panorama computazionale si divide in due famiglie concettualmente distinte: i \textbf{metodi diretti} e i \textbf{metodi iterativi}.

\begin{table}[htbp]
\centering
\small
\renewcommand{\arraystretch}{1.3}
\begin{tabularx}{\textwidth}{lYY}
\toprule
\textbf{Caratteristica} & \textbf{Metodi Diretti ($LU, QR$, Cholesky)} & \textbf{Metodi Iterativi (Krylov, CG, GMRES)} \\
\midrule
\textbf{Principio Base} & Fattorizzazione matriciale esatta di $A$ in matrici strutturate semplici. & Costruzione di una successione convergente di vettori $x^{(0)}, x^{(1)}, \dots \to x^*$. \\
\textbf{Accuratezza} & Soluzione esatta in aritmetica esatta (affetta solo da round-off in virgola mobile). & Soluzione approssimata entro una tolleranza prefissata $\|Ax^{(k)} - b\| \le \varepsilon$. \\
\textbf{Costo Matrici Dense} & Rigido: $\mathcal{O}(n^3)$ operazioni a monte, indipendentemente dal vettore $b$. & Flessibile: $k$ iterazioni, ciascuna a costo $\mathcal{O}(n^2)$ (prodotto matrice-vettore). \\
\textbf{Preservazione Sparsit\`a} & Spesso soggetta a \textit{fill-in} (generazione di elementi non nulli nei fattori $L$ e $U$). & Perfetta: la matrice $A$ \`e utilizzata esclusivamente come operatore \textit{black-box} $v \mapsto Av$. \\
\textbf{Taglia Tipica} & Matrici medio-piccole ($n \lesssim 10^4$ su workstation standard). & Matrici enormi o sparse derivanti da PDE ($n \approx 10^6 \div 10^9$). \\
\bottomrule
\end{tabularx}
\caption{Confronto tassonomico comparativo tra metodi diretti e metodi iterativi.}
\label{tab:diretti-vs-iterativi}
\end{table}

I metodi diretti eccellono quando:
\begin{enumerate}
    \item La dimensione della matrice $n$ consente di allocare i fattori in memoria RAM;
    \item \`E richiesta la massima accuratezza di macchina senza dipendere dalla velocit\`a di convergenza spettrale;
    \item Si devono risolvere \textbf{molteplici sistemi lineari} aventi la medesima matrice di coefficienti $A$ ma vettori termine noto $b^{(1)}, b^{(2)}, \dots, b^{(p)}$ differenti: la fattorizzazione $\mathcal{O}(n^3)$ viene eseguita una sola volta \textit{una tantum}, dopodich\'e ciascun sistema successivo costa solo $\mathcal{O}(n^2)$.
\end{enumerate}

\FloatBarrier
\section{La Fattorizzazione LU Standard}
\label{sec:fattorizzazione-lu-standard}

\subsection{Definizione e Schema di Risoluzione a Due Fasi}
\begin{definition}[Fattorizzazione LU]
Data una matrice quadrata $A \in \R^{n \times n}$, la sua \textbf{fattorizzazione $LU$} \`e la decomposizione:
\begin{equation}
    A = L U,
    \label{eq:decomposizione-lu}
\end{equation}
dove:
\begin{itemize}
    \item $L \in \R^{n \times n}$ \`e una matrice triangolare inferiore a \textbf{diagonale unitaria} (\textit{unit lower triangular}):
    \begin{equation}
        l_{ii} = 1 \quad (\forall i), \qquad l_{ij} = 0 \quad (\forall j > i);
    \end{equation}
    \item $U \in \R^{n \times n}$ \`e una matrice triangolare superiore (\textit{upper triangular}):
    \begin{equation}
        u_{ij} = 0 \quad (\forall i > j).
    \end{equation}
\end{itemize}
\end{definition}

Se la fattorizzazione $A = LU$ \`e nota, la risoluzione di $Ax = b$ si articola in due sottoproblemi triangolari a cascata ponendo la variabile ausiliaria $y := Ux$:
\begin{equation}
    A x = b \iff L \underbrace{(U x)}_{y} = b \iff
    \begin{cases}
        L y = b & \text{(Fase 1: Forward Substitution)} \\
        U x = y & \text{(Fase 2: Backward Substitution)}
    \end{cases}
    \label{eq:schema-due-fasi-lu}
\end{equation}

\begin{figure}[htbp]
\centering
\begin{tikzpicture}[
    >=Stealth,
    scale=0.81,
    every node/.style={transform shape},
    block/.style={
        rectangle, thick, rounded corners=4pt,
        minimum height=1.15cm, minimum width=4.7cm,
        align=center, font=\small
    },
    inout/.style={
        block, draw=navy!85, fill=blue!5
    },
    proc/.style={
        block, draw=navy!85, fill=navy!6
    },
    solvebox/.style={
        block, draw=teal!85!black, fill=teal!6
    },
    solbox/.style={
        block, draw=green!60!black, fill=green!7
    },
    lbl/.style={font=\scriptsize\bfseries, align=center, fill=white, inner sep=2.5pt},
    arr/.style={->, thick, draw=navy!85}
]

  % ========================================================
  % COLONNA SINISTRA: FASE 0 - FATTORIZZAZIONE (UNA TANTUM)
  % ========================================================
  \node[inout] (in_A)   at (0, 3.6) {\textbf{Matrice di Input}\\$A \in \R^{n \times n}$};
  \node[proc]  (fact)   at (0, 1.95) {\textbf{Eliminazione Gaussiana}\\Decomposizione $PA = LU$};
  \node[proc]  (out_LU) at (0, 0.3)  {\textbf{Fattori di Decomposizione}\\$L, U$ triangolari, $P$ permutazione};

  % ========================================================
  % COLONNA DESTRA: FASI DI SOLVE (PER CIASCUN TERMINE NOTO)
  % Allargato a X = 9.4 per garantire un ampio corridoio centrale
  % ========================================================
  \node[inout]    (in_b)  at (9.4, 3.6) {\textbf{Termine Noto di Input}\\$b \in \R^n$};
  \node[solvebox] (fwd)   at (9.4, 1.4) {\textbf{Fase 1: Forward Substitution}\\$Ly = \tilde{b} \quad (\tilde{b} = Pb)$};
  \node[solvebox] (bwd)   at (9.4, -0.8) {\textbf{Fase 2: Backward Substitution}\\$Ux = y$};
  \node[solbox]   (out_x) at (9.4, -2.8) {\textbf{Soluzione Finale $x \in \R^n$}\\$Ax = b$ \textbf{(senza calcolare $A^{-1}$)}};

  % ========================================================
  % CONTENITORI DI SFONDO (Tratteggiati)
  % ========================================================
  \begin{scope}[on background layer]
    % Container Fattorizzazione (Sinistra)
    \node[draw=navy!70, dashed, thick, rounded corners=6pt,
          fit=(in_A) (fact) (out_LU),
          inner xsep=12pt, inner ysep=16pt] (box_fact) {};
    \node[font=\footnotesize\bfseries, text=navy!85, fill=white, inner sep=3pt] 
          at (box_fact.north) 
          {Fattorizzazione ($\mathcal{O}(n^3)$)};

    % Container Risoluzione (Destra)
    \node[draw=teal!70!black, dashed, thick, rounded corners=6pt,
          fit=(in_b) (fwd) (bwd) (out_x),
          inner xsep=12pt, inner ysep=16pt] (box_solve) {};
    \node[font=\footnotesize\bfseries, text=teal!85!black, fill=white, inner sep=3pt] 
          at (box_solve.north) 
          {Risoluzione a Due Stadi ($\mathcal{O}(n^2)$)};
  \end{scope}

  % Connessioni verticali colonna sinistra
  \draw[arr] (in_A) -- (fact);
  \draw[arr] (fact) -- (out_LU);

  % Connessioni verticali colonna destra
  \draw[arr] (in_b) -- (fwd);
  \draw[arr] (fwd) -- node[midway, right=8pt, font=\small\bfseries, text=teal!85!black] {$y \in \R^n$} (bwd);
  \draw[arr] (bwd) -- (out_x);

  % Connessioni trasversali ortogonali simmetriche (bus centrale a X = 4.4)
  \coordinate (bus) at (4.4, 0.3);
  \coordinate (c_fwd) at (4.4, 1.4);
  \coordinate (c_bwd) at (4.4, -0.8);
  
  \draw[thick, draw=teal!80!black] (out_LU.east) -- (bus);
  \draw[thick, draw=teal!80!black] (bus) -- (c_fwd);
  \draw[thick, draw=teal!80!black] (bus) -- (c_bwd);
  
  % Etichette posizionate entrambe SOPRA (above=2pt) nel corridoio centrale prima del box tratteggiato
  \draw[arr, draw=teal!80!black] (c_fwd) -- 
      node[lbl, pos=0.34, above=2pt] {Fattore $L$ (e $P$)} (fwd.west);
      
  \draw[arr, draw=teal!80!black] (c_bwd) -- 
      node[lbl, pos=0.34, above=2pt] {Fattore $U$} (bwd.west);

\end{tikzpicture}
\caption{Pipeline computazionale completa per la risoluzione di un sistema lineare $Ax = b$ mediante eliminazione gaussiana e fattorizzazione $PA = LU$: netta separazione tra la decomposizione matriciale a monte ($\approx \frac{2}{3}n^3$ FLOPs, calcolata \textit{una tantum}) e la risoluzione triangolare a due stadi ($2n^2$ FLOPs, ripetibile economicamente per qualsiasi termine noto $b$). La soluzione finale $x$ viene ottenuta senza mai formare la matrice inversa $A^{-1}$.}
\label{fig:pipeline-lu}
\end{figure}

L'ordine sequenziale \`e vincolante: non \`e possibile determinare $x$ prima di aver ricavato $y$. Ciascuna delle due fasi richiede $\mathcal{O}(n^2)$ operazioni, per un costo di risoluzione di $2n^2 = \mathcal{O}(n^2)$ FLOPs.

\subsection{Derivazione Costruttiva: Eliminazione Gaussiana e Matrici Elementari}
Come si calcolano concretamente i fattori $L$ e $U$? L'algoritmo coincide con l'\textbf{eliminazione gaussiana}, formalizzata algebricamente attraverso il prodotto di matrici elementari di trasformazione.

\subsubsection{Passo 1 dell'Eliminazione Gaussiana}
Data $A \in \R^{n \times n}$, supponiamo che l'elemento pivot $a_{11} \neq 0$. Per annullare gli elementi sottodiagonali della prima colonna ($a_{21}, a_{31}, \dots, a_{n1}$), per ogni riga $i = 2, \dots, n$ si definisce il moltiplicatore:
\begin{equation}
    m_{i1} := \frac{a_{i1}}{a_{11}},
\end{equation}
e si sottrae alla riga $i$ il multiplo $m_{i1}$ della riga 1:
\begin{equation}
    R_i \leftarrow R_i - m_{i1} R_1.
\end{equation}

Tale operazione equivale a premoltiplicare la matrice $A$ per la \textbf{matrice elementare di Gauss} $E_1 \in \R^{n \times n}$:
\begin{equation}
    E_1 = \begin{bmatrix}
        1 & 0 & 0 & \dots & 0 \\
        -m_{21} & 1 & 0 & \dots & 0 \\
        -m_{31} & 0 & 1 & \dots & 0 \\
        \vdots & \vdots & \vdots & \ddots & \vdots \\
        -m_{n1} & 0 & 0 & \dots & 1
    \end{bmatrix} = I_n - m^{(1)} e_1^T,
    \label{eq:matrice-elementare-e1}
\end{equation}
dove $m^{(1)} = [0, m_{21}, m_{31}, \dots, m_{n1}]^T$ ed $e_1 = [1, 0, \dots, 0]^T$. La matrice risultante dopo il primo passo \`e:
\begin{equation}
    A_1 = E_1 A = \begin{bmatrix}
        a_{11} & a_{12} & \dots & a_{1n} \\
        0 & a_{22}^{(1)} & \dots & a_{2n}^{(1)} \\
        \vdots & \vdots & \ddots & \vdots \\
        0 & a_{n2}^{(1)} & \dots & a_{nn}^{(1)}
    \end{bmatrix}.
\end{equation}

\subsubsection{Passi Successivi e Ricostruzione di $U$}
Iterando il procedimento per $k = 1, 2, \dots, n-1$, assumendo che a ogni stadio il pivot $a_{kk}^{(k-1)} \neq 0$, si azzerano gli elementi della colonna $k$ al di sotto della diagonale mediante la matrice:
\begin{equation}
    E_k = I_n - m^{(k)} e_k^T, \qquad \text{con } m^{(k)} = [\underbrace{0, \dots, 0}_{k \text{ zeri}}, m_{k+1, k}, \dots, m_{nk}]^T,
\end{equation}
in cui $m_{ik} = \frac{a_{ik}^{(k-1)}}{a_{kk}^{(k-1)}}$. Al termine di $n-1$ passaggi, la matrice risulta triangolare superiore e coincide con $U$:
\begin{equation}
    U = E_{n-1} E_{n-2} \dots E_2 E_1 A.
    \label{eq:relazione-u-prodotto-e}
\end{equation}

\subsubsection{Propriet\`a Inversa e Costruzione di $L$}
Invertendo la relazione~\eqref{eq:relazione-u-prodotto-e}, otteniamo:
\begin{equation}
    A = (E_{n-1} \dots E_2 E_1)^{-1} U = E_1^{-1} E_2^{-1} \dots E_{n-1}^{-1} U.
\end{equation}
Ponendo per definizione $L := E_1^{-1} E_2^{-1} \dots E_{n-1}^{-1}$, si ha $A = LU$.

\begin{property}[Inversa e Prodotto di Matrici Elementari di Gauss]
Le matrici elementari $E_k$ godono di due propriet\`a algebriche fondamentali:
\begin{enumerate}
    \item L'inversa di $E_k = I_n - m^{(k)} e_k^T$ si ottiene semplicemente \textbf{invertendo il segno dei moltiplicatori}:
    \begin{equation}
        E_k^{-1} = I_n + m^{(k)} e_k^T = \begin{bmatrix}
            1 & & & & \\
            & \ddots & & & \\
            & & 1 & & \\
            & & m_{k+1, k} & \ddots & \\
            & & m_{nk} & & 1
        \end{bmatrix}.
    \end{equation}
    \item Il prodotto ordinato $L = E_1^{-1} E_2^{-1} \dots E_{n-1}^{-1}$ non richiede alcun calcolo algebrico esplicito: le colonne extradiagonali dei vettori $m^{(k)}$ si \textbf{compenetrano perfettamente senza interferire}:
    \begin{equation}
        L = \begin{bmatrix}
            1 & 0 & 0 & \dots & 0 \\
            m_{21} & 1 & 0 & \dots & 0 \\
            m_{31} & m_{32} & 1 & \dots & 0 \\
            \vdots & \vdots & \vdots & \ddots & \vdots \\
            m_{n1} & m_{n2} & m_{n3} & \dots & 1
        \end{bmatrix}.
    \end{equation}
\end{enumerate}
\end{property}
La matrice $L$ possiede dunque elementi unitari sulla diagonale principale e contiene, nelle rispettive posizioni strettamente triangolari inferiori, \textbf{esattamente i moltiplicatori di Gauss} impiegati durante l'eliminazione.

\subsection{Esempio Numerico Dettagliato a Mano}
\label{subsec:esempio-numerico-lu-manuale}
Consideriamo la matrice $3 \times 3$ analizzata nella lezione:
\begin{equation}
    A = \begin{bmatrix}
        2 & 1 & -1 \\
        1 & 0 & 2 \\
        -4 & 3 & 2
    \end{bmatrix}.
\end{equation}

\paragraph{Passo 1 dell'Eliminazione:}
\begin{itemize}
    \item Pivot designato: $a_{11} = 2 \neq 0$.
    \item Moltiplicatore riga 2: $m_{21} = \frac{a_{21}}{a_{11}} = \frac{1}{2}$.
    \item Moltiplicatore riga 3: $m_{31} = \frac{a_{31}}{a_{11}} = \frac{-4}{2} = -2$.
    \item Aggiornamento delle righe:
    \begin{align}
        R_2 &\leftarrow R_2 - \tfrac{1}{2} R_1 = \left[0, \; -\tfrac{1}{2}, \; \tfrac{5}{2}\right], \\
        R_3 &\leftarrow R_3 - (-2) R_1 = [0, \; 5, \; 0].
    \end{align}
\end{itemize}
Matrice al termine del passo 1 e corrispondente matrice elementare $E_1$:
\begin{equation}
    A_1 = \begin{bmatrix}
        2 & 1 & -1 \\
        0 & -\frac{1}{2} & \frac{5}{2} \\
        0 & 5 & 0
    \end{bmatrix}, \qquad
    E_1 = \begin{bmatrix}
        1 & 0 & 0 \\
        -\frac{1}{2} & 1 & 0 \\
        2 & 0 & 1
    \end{bmatrix}.
\end{equation}

\paragraph{Passo 2 dell'Eliminazione:}
\begin{itemize}
    \item Nuovo pivot: $a_{22}^{(1)} = -\frac{1}{2} \neq 0$.
    \item Moltiplicatore riga 3: $m_{32} = \frac{a_{32}^{(1)}}{a_{22}^{(1)}} = \frac{5}{-1/2} = -10$.
    \item Aggiornamento riga 3:
    \begin{equation}
        R_3 \leftarrow R_3 - (-10) R_2 = \left[0, \; 5 - 10\left(\frac{1}{2}\right), \; 0 + 10\left(\frac{5}{2}\right)\right] = [0, \; 0, \; 25].
    \end{equation}
\end{itemize}
Matrice finale $U$ e matrice elementare $E_2$:
\begin{equation}
    U = \begin{bmatrix}
        2 & 1 & -1 \\
        0 & -\frac{1}{2} & \frac{5}{2} \\
        0 & 0 & 25
    \end{bmatrix}, \qquad
    E_2 = \begin{bmatrix}
        1 & 0 & 0 \\
        0 & 1 & 0 \\
        0 & 10 & 1
    \end{bmatrix}.
\end{equation}

\paragraph{Ricostruzione del Fattore $L$:}
Invertendo i segni sottodiagonali di $E_1$ ed $E_2$:
\begin{equation}
    E_1^{-1} = \begin{bmatrix}
        1 & 0 & 0 \\
        \frac{1}{2} & 1 & 0 \\
        -2 & 0 & 1
    \end{bmatrix}, \qquad
    E_2^{-1} = \begin{bmatrix}
        1 & 0 & 0 \\
        0 & 1 & 0 \\
        0 & -10 & 1
    \end{bmatrix}.
\end{equation}
Moltiplicando le due matrici:
\begin{equation}
    L = E_1^{-1} E_2^{-1} = \begin{bmatrix}
        1 & 0 & 0 \\
        \frac{1}{2} & 1 & 0 \\
        -2 & -10 & 1
    \end{bmatrix}.
\end{equation}
Verifica del prodotto:
\begin{align}
    LU &= \begin{bmatrix}
        1 & 0 & 0 \\
        \frac{1}{2} & 1 & 0 \\
        -2 & -10 & 1
    \end{bmatrix}
    \begin{bmatrix}
        2 & 1 & -1 \\
        0 & -\frac{1}{2} & \frac{5}{2} \\
        0 & 0 & 25
    \end{bmatrix} \nonumber \\
    &= \begin{bmatrix}
        2 & 1 & -1 \\
        1 & 0 & 2 \\
        -4 & 3 & 2
    \end{bmatrix} = A.
\end{align}
La fattorizzazione $LU$ calcolata \`e esatta.

\subsection{Complessit\`a Computazionale della Fattorizzazione $LU$}
Determiniamo il conteggio esatto delle operazioni floating-point per una matrice $A \in \R^{n \times n}$.
Al generico passo $k \in \{1, \dots, n-1\}$:
\begin{enumerate}
    \item Si calcolano gli $n - k$ moltiplicatori $m_{ik} = a_{ik}/a_{kk}$, richiedendo $n - k$ divisioni;
    \item Si aggiornano gli elementi sottomatriciali: per ciascuna delle $n - k$ righe sottostanti, si aggiornano le $n - k$ colonne, comportando $(n - k)^2$ moltiplicazioni e $(n - k)^2$ sottrazioni.
\end{enumerate}
Sommando su tutti gli stadi $k = 1, \dots, n-1$:
\begin{equation}
\begin{aligned}
    \sum_{k=1}^{n-1} (n - k)^2 &= \sum_{j=1}^{n-1} j^2 = \frac{(n-1)n(2n-1)}{6} \\
    &= \frac{2n^3 - 3n^2 + n}{6} \approx \frac{n^3}{3} \text{ prodotti e } \frac{n^3}{3} \text{ somme.}
\end{aligned}
\end{equation}
Tenendo conto anche delle divisioni:
\begin{equation}
    \text{FLOPs Totali} \approx 2 \cdot \frac{n^3}{3} = \frac{2}{3} n^3 \implies \mathcal{O}(n^3).
\end{equation}

\begin{property}[Conservazione della Memoria: Algoritmo In-Place]
L'eliminazione gaussiana non richiede l'allocazione di nuove matrici in memoria. Man mano che un elemento sottodiagonale $a_{ik}$ viene azzerato dall'algoritmo, la sua posizione fisica nella matrice in memoria viene immediatamente riutilizzata per memorizzare il rispettivo moltiplicatore $l_{ik} = m_{ik}$. La diagonale unitaria di $L$ ($l_{ii} = 1$) non viene memorizzata. Al termine dell'esecuzione, la matrice di input $A$ contiene sia $L$ (nella parte strettamente triangolare inferiore) sia $U$ (nella parte triangolare superiore e sulla diagonale principale).
\end{property}

\FloatBarrier
\section{Fallimento di LU, Instabilit\`a Numerica e Pivoting Parziale}
\label{sec:fallimento-instabilita-pivoting}

L'algoritmo di eliminazione gaussiana formulato nella Sezione~\ref{sec:fattorizzazione-lu-standard} soffre di due vulnerabilit\`a intrinseche:
\begin{enumerate}
    \item \textbf{Fallimento Algebrico per Pivot Nullo:} Se l'elemento diagonale al passo $k$ risulta esattamente nullo ($a_{kk}^{(k-1)} = 0$), l'algoritmo fallisce immediatamente a causa di una divisione per zero, persino se la matrice $A$ \`e perfettamente invertibile;
    \item \textbf{Catastrofe Numerica per Pivot Piccolo:} Se il pivot \`e non nullo ma di modulo estremamente ridotto rispetto agli altri elementi della sottomatrice, i moltiplicatori di Gauss esplodono, provocando gravissima perdita di precisione dovuta a cancellazione numerica ed errori di arrotondamento in aritmetica a virgola mobile.
\end{enumerate}

\subsection{Problema 1: Pivot Esattamente Nullo}
Si consideri la matrice non singolare:
\begin{equation}
    A = \begin{bmatrix}
        0 & 1 & 2 \\
        3 & 4 & 5 \\
        6 & 7 & 8
    \end{bmatrix}.
\end{equation}
Poich\'e $a_{11} = 0$, \`e impossibile calcolare $m_{21} = 3/0$ e $m_{31} = 6/0$: l'eliminazione gaussiana standard collassa al primo passo.
Tuttavia, poich\'e le righe di un sistema lineare rappresentano singole equazioni, l'ordine con cui esse sono elencate \`e del tutto arbitrario. Scambiando la prima riga con la seconda, si ottiene il sistema equivalente:
\begin{equation}
    \begin{bmatrix}
        3 & 4 & 5 \\
        0 & 1 & 2 \\
        6 & 7 & 8
    \end{bmatrix}
    \begin{bmatrix} x_1 \\ x_2 \\ x_3 \end{bmatrix}
    =
    \begin{bmatrix} b_2 \\ b_1 \\ b_3 \end{bmatrix},
\end{equation}
il cui pivot iniziale \`e $3 \neq 0$, consentendo all'algoritmo di proseguire regolarmente.

\subsection{Problema 2: Instabilit\`a da Pivot Piccolo (Il Controesempio Patologico)}
La vulnerabilit\`a pi\`u insidiosa emerge quando il pivot non \`e zero, ma un numero piccolo in valore assoluto.

\begin{example}[Instabilit\`a dell'Eliminazione Ingenua]
\label{ex:bad-example-pivot}
Si consideri la matrice:
\begin{equation}
    A = \begin{bmatrix}
        10^{-10} & 1 & 2 \\
        1 & 4 & 5 \\
        -1 & 7 & 8
    \end{bmatrix}.
\end{equation}
Al passo 1, il pivot \`e $a_{11} = 10^{-10} \neq 0$. Eseguiamo il passo di eliminazione:
\begin{itemize}
    \item Moltiplicatore riga 2: $m_{21} = \frac{1}{10^{-10}} = 10^{10}$;
    \item Moltiplicatore riga 3: $m_{31} = \frac{-1}{10^{-10}} = -10^{10}$;
    \item Aggiornamento della riga 2:
    \begin{equation}
        R_2 \leftarrow R_2 - 10^{10} R_1 = \left[0, \; 4 - 10^{10}, \; 5 - 2 \cdot 10^{10}\right].
    \end{equation}
\end{itemize}
In aritmetica floating-point a precisione standard (IEEE 754 a 64 bit, con precisione di macchina $\varepsilon_{\mathrm{mach}} \approx 10^{-16}$), la presenza del termine dominante $10^{10}$ cancella completamente le cifre dei coefficienti originari:
\begin{equation}
    \mathrm{fl}(4 - 10^{10}) = -10^{10}, \qquad \mathrm{fl}(5 - 2 \cdot 10^{10}) = -2 \cdot 10^{10}.
\end{equation}
Il nuovo pivot al passo 2 avr\`a modulo dell'ordine di $10^{10}$. Di conseguenza, il fattore $U$ conterr\`a contemporaneamente elementi di grandezza originaria ($10^{-10}$) ed elementi artificialmente giganteschi ($10^{10}$).
\end{example}

\subsubsection{Analisi di Stabilit\`a all'Indietro (Backward Error Analysis)}
Per comprendere teoricamente la gravit\`a del fenomeno, richiamiamo il teorema fondamentale di backward stability per la fattorizzazione $LU$ calcolata in floating-point~\cite{golub2013matrix,higham2002accuracy}:
\begin{theorem}[Backward Error della Fattorizzazione LU]
Siano $\tilde{L}$ e $\tilde{U}$ i fattori calcolati numericamente in aritmetica a precisione $\varepsilon_{\mathrm{mach}}$. Essi rappresentano la fattorizzazione $LU$ esatta di una matrice perturbata $\tilde{A} = \tilde{L}\tilde{U} = A + \Delta A$, per la quale vale la maggiorazione:
\begin{equation}
    \|\Delta A\| = \|A - \tilde{L}\tilde{U}\| \le \gamma_n \|\tilde{L}\| \|\tilde{U}\| \varepsilon_{\mathrm{mach}},
    \label{eq:backward-error-lu}
\end{equation}
dove $\gamma_n = \frac{c \cdot n \varepsilon_{\mathrm{mach}}}{1 - c \cdot n \varepsilon_{\mathrm{mach}}} = \mathcal{O}(n \varepsilon_{\mathrm{mach}})$.
\end{theorem}

L'errore di retroazione $\|\Delta A\|$ \`e proporzionale al prodotto delle norme dei fattori $\|\tilde{L}\| \|\tilde{U}\|$.
\begin{itemize}
    \item Se gli elementi di $\tilde{L}$ e $\tilde{U}$ rimangono comparabili a quelli di $A$, si ha $\|\tilde{L}\| \|\tilde{U}\| \approx \|A\|$, garantendo che $\|\Delta A\| / \|A\| \approx \mathcal{O}(\varepsilon_{\mathrm{mach}})$ (algoritmo backward stable);
    \item Se, come nell'Esempio~\ref{ex:bad-example-pivot}, gli elementi esplodono a $\|\tilde{U}\| \approx 10^{10}$, allora:
    \begin{equation}
        \|\Delta A\| \approx 10^{10} \cdot 10^{-16} = 10^{-6}.
    \end{equation}
    Si sono irrimediabilmente perdute \textbf{10 cifre decimali di precisione}: l'algoritmo \`e numericamente instabile e la soluzione computata sar\`a inaffidabile.
\end{itemize}

\subsection{Strategia del Pivoting Parziale (Partial Pivoting)}
\label{subsec:pivoting-parziale}

Per neutralizzare sia le divisioni per zero sia l'esplosione dei moltiplicatori, si applica la tecnica del \textbf{pivoting parziale}:
\begin{definition}[Regola del Pivoting Parziale]
Al generico stadio $k \in \{1, \dots, n-1\}$, tra tutti gli elementi della colonna corrente situati sulla diagonale o al di sotto di essa (ovvero lungo le righe $i = k, k+1, \dots, n$), si ricerca l'elemento avente \textbf{massimo valore assoluto}:
\begin{equation}
    |a_{rk}^{(k-1)}| = \max_{k \le i \le n} |a_{ik}^{(k-1)}|.
    \label{eq:regola-pivoting-parziale}
\end{equation}
Si scambia quindi fisicamente la riga corrente $k$ con la riga $r$.
\end{definition}

Effettuando questo scambio prima di calcolare i moltiplicatori, per ogni riga sottostante $i > k$ si garantisce che:
\begin{equation}
    |m_{ik}| = \left| \frac{a_{ik}^{(k-1)}}{a_{kk}^{(k-1)}} \right| \le 1, \qquad \forall i = k+1, \dots, n.
\end{equation}
Tutti i moltiplicatori di Gauss (e quindi tutti gli elementi extradiagonali della matrice $L$) hanno valore assoluto strettamente non superiore a 1:
\begin{equation}
    |l_{ij}| \le 1 \quad (\forall i > j) \implies \|L\|_{\infty} \le n.
\end{equation}
Ci\`o previene l'esplosione della norma $\|L\|$ e stabilizza incondizionatamente il calcolo numerico.

\subsection{Formalizzazione con Matrici di Permutazione: Teorema $PA = LU$}
Lo scambio di righe si formalizza algebricamente mediante le \textbf{matrici di permutazione}.

\begin{definition}[Matrice di Permutazione]
Una matrice $P \in \R^{n \times n}$ si dice matrice di permutazione se \`e ottenuta permutando le righe (o le colonne) della matrice identit\`a $I_n$.
\end{definition}

\begin{property}[Propriet\`a delle Matrici di Permutazione]
Le matrici di permutazione soddisfano:
\begin{enumerate}
    \item $P$ \`e ortogonale: $P^T P = P P^T = I_n \implies P^{-1} = P^T$;
    \item Moltiplicare a sinistra $PA$ permuta le righe di $A$ secondo la permutazione definita da $P$;
    \item Moltiplicare a destra $AP$ permuta le colonne di $A$.
\end{enumerate}
\end{property}

\begin{theorem}[Fattorizzazione $PA = LU$]
\label{thm:pa-lu}
Per ogni matrice quadrata non singolare $A \in \R^{n \times n}$, esistono una matrice di permutazione $P$, una matrice triangolare inferiore a diagonale unitaria $L$ (con $|l_{ij}| \le 1$ per ogni $i > j$) e una matrice triangolare superiore non singolare $U$ tali che:
\begin{equation}
    P A = L U.
    \label{eq:teorema-pa-lu}
\end{equation}
\end{theorem}

\subsubsection{Risoluzione del Sistema con $PA = LU$}
Dato il sistema originario $Ax = b$, premoltiplicando per la matrice di permutazione $P$:
\begin{equation}
    P A x = P b.
\end{equation}
Poich\'e $PA = LU$, si ottiene l'equazione $LUx = Pb$, che si risolve con lo schema a due stadi:
\begin{enumerate}
    \item \textbf{Permutazione del termine noto e Forward Substitution:}
    \begin{equation}
        L y = P b;
    \end{equation}
    \item \textbf{Backward Substitution:}
    \begin{equation}
        U x = y.
    \end{equation}
\end{enumerate}
Il calcolo di $Pb$ non richiede moltiplicazioni aritmetiche: consiste esclusivamente nel riordinare le componenti del vettore $b$ secondo il vettore degli indici di permutazione.

\subsection{Esempio Numerico Rivisitato con Pivoting Parziale}
Riprendiamo la matrice analizzata nella Sezione~\ref{subsec:esempio-numerico-lu-manuale}:
\begin{equation}
    A = \begin{bmatrix}
        2 & 1 & -1 \\
        1 & 0 & 2 \\
        -4 & 3 & 2
    \end{bmatrix}.
\end{equation}

\paragraph{Passo 1 (Ricerca del Pivot Massimo e Scambio):}
I candidati della prima colonna sono gli elementi di modulo $|2| = 2$, $|1| = 1$ e $|-4| = 4$.
Il massimo assoluto \`e $4$, situato nella riga 3.
Scambiamo la riga 1 con la riga 3 (e successivamente ordiniamo le righe in modo che la riga 3 diventi la prima riga, la riga 1 diventi la seconda e la riga 2 diventi la terza):
\begin{equation}
    PA = \begin{bmatrix}
        -4 & 3 & 2 \\
        2 & 1 & -1 \\
        1 & 0 & 2
    \end{bmatrix}, \qquad
    P = \begin{bmatrix}
        0 & 0 & 1 \\
        1 & 0 & 0 \\
        0 & 1 & 0
    \end{bmatrix}.
\end{equation}

\paragraph{Eliminazione al Passo 1:}
\begin{itemize}
    \item Pivot: $-4$;
    \item Moltiplicatore riga 2: $m_{21} = \frac{2}{-4} = -\frac{1}{2}$ (notiamo $|-1/2| \le 1$);
    \item Moltiplicatore riga 3: $m_{31} = \frac{1}{-4} = -\frac{1}{4}$ (notiamo $|-1/4| \le 1$);
    \item Aggiornamenti:
    \begin{align}
        R_2 &\leftarrow R_2 + \tfrac{1}{2} R_1 = \left[0, \; \tfrac{5}{2}, \; 0\right], \\
        R_3 &\leftarrow R_3 + \tfrac{1}{4} R_1 = \left[0, \; \tfrac{3}{4}, \; \tfrac{5}{2}\right].
    \end{align}
\end{itemize}

\paragraph{Passo 2 (Verifica del Pivot e Conclusione):}
Nella seconda colonna, i candidati sono gli elementi sottodiagonali: in riga 2 abbiamo $\frac{5}{2} = 2.5$, in riga 3 abbiamo $\frac{3}{4} = 0.75$.
Poich\'e $\frac{5}{2} > \frac{3}{4}$, il massimo si trova gi\`a sulla diagonale in riga 2: \textbf{non occorre alcuno scambio di righe}.
\begin{itemize}
    \item Moltiplicatore riga 3: $m_{32} = \frac{3/4}{5/2} = \frac{3}{4} \cdot \frac{2}{5} = \frac{3}{10}$ (notiamo $|3/10| \le 1$);
    \item Aggiornamento riga 3:
    \begin{equation}
        R_3 \leftarrow R_3 - \frac{3}{10} R_2 = \left[0, \; \frac{3}{4} - \frac{3}{10}\left(\frac{5}{2}\right), \; \frac{5}{2} - \frac{3}{10}(0)\right] = \left[0, \; 0, \; \frac{5}{2}\right].
    \end{equation}
\end{itemize}

\paragraph{Fattori $P, L, U$ Risultanti:}
\begin{equation}
    P = \begin{bmatrix}
        0 & 0 & 1 \\
        1 & 0 & 0 \\
        0 & 1 & 0
    \end{bmatrix}, \qquad
    L = \begin{bmatrix}
        1 & 0 & 0 \\
        -\frac{1}{2} & 1 & 0 \\
        -\frac{1}{4} & \frac{3}{10} & 1
    \end{bmatrix}, \qquad
    U = \begin{bmatrix}
        -4 & 3 & 2 \\
        0 & \frac{5}{2} & 0 \\
        0 & 0 & \frac{5}{2}
    \end{bmatrix}.
\end{equation}
Tutti i moltiplicatori in $L$ soddisfano $|l_{ij}| \le 1$, garantendo stabilit\`a numerica ottimale.

\subsection{Costo del Pivoting e Cenni sul Pivoting Totale}
\begin{property}[Costo Asintotico del Pivoting Parziale]
A ogni passo $k$, individuare il pivot massimo richiede $n - k$ confronti scalari tra numeri reali. Sommando su tutti gli stadi:
\begin{equation}
    \sum_{k=1}^{n-1} (n - k) = \frac{n(n-1)}{2} = \mathcal{O}(n^2) \text{ confronti.}
\end{equation}
Poich\'e la ricerca del pivot e gli scambi logici di puntatori hanno costo quadratico $\mathcal{O}(n^2)$, tale overhead \`e \textbf{completamente trascurabile} rispetto alle $\frac{2}{3}n^3$ operazioni aritmetiche dell'eliminazione gaussiana.
\end{property}

\begin{remark}[Pivoting Totale e Caso Patologico di Wilkinson]
In teoria, esistono rare matrici patologiche (costruite artificialmente da J. H. Wilkinson) per le quali il fattore di crescita degli elementi di $U$ con pivoting parziale pu\`o raggiungere nel caso peggiore $2^{n-1}$~\cite{trefethen1997numerical}. Per prevenire anche questo scenario teorico estremo, si potrebbe adottare il \textbf{pivoting totale} (\textit{Complete Pivoting}), ricercando il massimo su tutta la sottomatrice $(n-k+1) \times (n-k+1)$ e scambiando sia righe che colonne:
\begin{equation}
    P A Q = L U,
\end{equation}
dove $Q$ \`e una matrice di permutazione delle colonne. Tuttavia, il pivoting totale richiede $\mathcal{O}(n^3)$ confronti, raddoppiando o triplicando il tempo di calcolo. Poich\'e nelle applicazioni reali la crescita $2^{n-1}$ non si riscontra praticamente mai, il \textbf{pivoting parziale costituisce lo standard de facto industriale e scientifico}.
\end{remark}

\FloatBarrier
\section{Matrici Singolari e Sistemi Rettangolari}
\label{sec:matrici-singolari-rettangolari}

Cosa accade alla fattorizzazione $LU$ se la matrice $A$ \`e singolare o rettangolare?
\begin{enumerate}
    \item \textbf{Esistenza della Forma a Scalini:} L'algoritmo di eliminazione gaussiana continua a essere applicabile; in presenza di dipendenza lineare, la matrice $U$ conterr\`a uno o pi\`u zeri sulla diagonale principale o intere righe nulle al fondo;
    \item \textbf{Caratterizzazione dello Spazio delle Colonne e del Nucleo:} La forma ridotta consente di identificare le colonne pivot per estrarre una base ortonormale di $\mathrm{Range}(A)$ e determinare una base del nucleo $\ker(A)$;
    \item \textbf{Sistemi Sovradeterminati o Sottodeterminati:} Se il fine applicativo \`e risolvere un sistema rettangolare $A \in \R^{m \times n}$ (tipicamente sovradeterminato nei problemi di regressione con $m \gg n$), la fattorizzazione $LU$ non \`e il metodo d'elezione. In tali casi, il sistema si riformula come problema ai \textbf{minimi quadrati lineari} (\textit{Ordinary Least Squares}):
    \begin{equation}
        \min_{x \in \R^n} \|Ax - b\|_2^2,
    \end{equation}
    risolto in modo numericamente superiore mediante la fattorizzazione $QR$ o la decomposizione ai valori singolari (SVD).
\end{enumerate}

\FloatBarrier
\section{Risoluzione di Sistemi Lineari con Fattorizzazioni Alternative: QR e SVD}
\label{sec:risoluzione-qr-svd}

Nelle lezioni precedenti sono state analizzate due fondamentali decomposizioni matriciali: la fattorizzazione $QR$ (Capitolo~\ref{chap:fattorizzazione-qr-riflessioni-householder}) e la decomposizione ai valori singolari SVD (Capitoli~\ref{chap:decomposizione-valori-singolari-svd} e~\ref{chap:svd-norme-matriciali-applicazioni}).
Se tali decomposizioni sono gi\`a disponibili a monte (o calcolate per scopi di analisi dati), come si risolve in modo efficiente il sistema $Ax = b$?

\subsection{Esercizio 1: Risoluzione mediante Fattorizzazione $QR$}
Sia nota la fattorizzazione $QR$ di una matrice quadrata invertibile $A \in \R^{n \times n}$:
\begin{equation}
    A = Q R,
\end{equation}
con $Q \in \R^{n \times n}$ ortogonale ($Q^T Q = I_n$) e $R \in \R^{n \times n}$ triangolare superiore non singolare.

Sostituendo nel sistema originario:
\begin{equation}
    Q R x = b.
\end{equation}
Moltiplicando a sinistra per la trasposta $Q^T$:
\begin{equation}
    Q^T Q R x = Q^T b \implies I_n R x = Q^T b \implies R x = Q^T b.
\end{equation}

\subsubsection{Algoritmo in Due Passi e Complessit\`a}
\begin{enumerate}
    \item \textbf{Proiezione Ortogonale del Termine Noto:} Si calcola il vettore proiettato:
    \begin{equation}
        c := Q^T b.
    \end{equation}
    Trattandosi di un prodotto matrice-vettore, il costo \`e di $2n^2 = \mathcal{O}(n^2)$ FLOPs;
    \item \textbf{Backward Substitution:} Si risolve il sistema triangolare superiore:
    \begin{equation}
        R x = c.
    \end{equation}
    Come dimostrato nella Sezione~\ref{sec:easy-cases}, l'algoritmo di sostituzione all'indietro richiede esattamente $n^2 = \mathcal{O}(n^2)$ FLOPs.
\end{enumerate}

\begin{property}[Costo di Risoluzione con $QR$ Precomputata]
Avendo gi\`a a disposizione i fattori $Q$ ed $R$, il costo computazionale totale per risolvere il sistema lineare \`e:
\begin{equation}
    \text{FLOPs} = 2n^2 + n^2 = 3n^2 \implies \mathcal{O}(n^2).
\end{equation}
\end{property}
Si noti che il calcolo \textit{ex novo} della fattorizzazione $QR$ mediante riflettori di Householder richiede $\frac{4}{3}n^3$ FLOPs, ovvero il doppio rispetto a $LU$ ($\frac{2}{3}n^3$). Pertanto, per matrici quadrate ordinarie si preferisce $LU$; $QR$ \`e invece imbattibile nei problemi ai minimi quadrati su sistemi rettangolari.

\subsection{Esercizio 2: Risoluzione mediante Decomposizione SVD}
Sia nota la decomposizione a valori singolari di una matrice non singolare $A \in \R^{n \times n}$:
\begin{equation}
    A = U \Sigma V^T,
\end{equation}
dove $U \in \R^{n \times n}$ e $V \in \R^{n \times n}$ sono matrici ortogonali ($U^T U = I_n$, $V^T V = I_n$) e $\Sigma = \diag(\sigma_1, \sigma_2, \dots, \sigma_n)$ con valori singolari strettamente positivi $\sigma_1 \ge \sigma_2 \ge \dots \ge \sigma_n > 0$.

Il sistema da risolvere \`e:
\begin{equation}
    U \Sigma V^T x = b.
\end{equation}
Ricordando la propriet\`a inversa del prodotto di matrici invertibili:
\begin{equation}
    (U \Sigma V^T)^{-1} = (V^T)^{-1} \Sigma^{-1} U^{-1}.
\end{equation}
Poich\'e $U$ e $V$ sono ortogonali, si ha $U^{-1} = U^T$ e $(V^T)^{-1} = (V^T)^T = V$. Dunque:
\begin{equation}
    x = V \Sigma^{-1} U^T b.
    \label{eq:soluzione-esplicita-svd}
\end{equation}

\subsubsection{Algoritmo Ottimizzato in Tre Passi}
L'espressione~\eqref{eq:soluzione-esplicita-svd} non va calcolata moltiplicando preliminarmente le tre matrici tra loro (che costerebbe $\mathcal{O}(n^3)$), ma associando da destra verso sinistra tramite prodotti matrice-vettore successivi:
\begin{enumerate}
    \item \textbf{Passo 1: Proiezione su $U$:} Si calcola il vettore ausiliario:
    \begin{equation}
        z := U^T b \in \R^n.
    \end{equation}
    Costo: prodotto matrice-vettore $\implies 2n^2 = \mathcal{O}(n^2)$ FLOPs;
    \item \textbf{Passo 2: Scalatura Diagonale con $\Sigma^{-1}$:} Poich\'e $\Sigma$ \`e diagonale, la sua inversa \`e $\Sigma^{-1} = \diag(1/\sigma_1, \dots, 1/\sigma_n)$. Il vettore $w := \Sigma^{-1} z$ si calcola componente per componente:
    \begin{equation}
        w_i = \frac{z_i}{\sigma_i}, \qquad \forall i = 1, \dots, n.
    \end{equation}
    Costo: esattamente $n$ divisioni scalari $\implies \mathcal{O}(n)$ FLOPs;
    \item \textbf{Passo 3: Ricostruzione tramite $V$:} Si ottiene la soluzione finale:
    \begin{equation}
        x := V w \in \R^n.
    \end{equation}
    Costo: prodotto matrice-vettore $\implies 2n^2 = \mathcal{O}(n^2)$ FLOPs.
\end{enumerate}

\begin{property}[Costo di Risoluzione con SVD Precomputata]
Disponendo della decomposizione SVD, la risoluzione esatta richiede:
\begin{equation}
    \text{FLOPs} = 2n^2 + n + 2n^2 = 4n^2 + n \implies \mathcal{O}(n^2).
\end{equation}
\end{property}

Il vantaggio supremo della risoluzione tramite SVD risiede nella capacit\`a di diagnosticare e trattare all'istante il \textbf{malcondizionamento numerico}: se alcuni valori singolari $\sigma_i$ sono inferiori alla precisione di macchina, \`e possibile troncarli ponendo $1/\sigma_i = 0$, ottenendo istantaneamente la soluzione regolarizzata a norma minima (pseudoinversa di Moore-Penrose $A^\dagger$).

\subsection{Sguardo al Futuro: Matrici Simmetriche e Fattorizzazione di Cholesky}
Se la matrice dei coefficienti $A$ possiede propriet\`a strutturali speciali, la complessit\`a della fattorizzazione diretta pu\`o essere ulteriormente abbattuta:
\begin{itemize}
    \item Se $A$ \`e \textbf{simmetrica} ($A = A^T$), l'eliminazione gaussiana preserva la simmetria, conducendo alla fattorizzazione:
    \begin{equation}
        A = L D L^T,
    \end{equation}
    con $L$ triangolare inferiore unitaria e $D$ diagonale, dimezzando sia la memoria che il tempo di calcolo a $\frac{1}{3}n^3$ FLOPs;
    \item Se $A$ \`e \textbf{simmetrica e definita positiva} (SPD, $x^T A x > 0$ per ogni $x \neq 0$), tutti i pivot diagonali sono garantiti strettamente positivi ($d_{ii} > 0$), rendendo il pivoting completamente superfluo. Definendo $L_{\mathrm{chol}} = L D^{1/2}$, si perviene alla celebre \textbf{Fattorizzazione di Cholesky}:
    \begin{equation}
        A = L_{\mathrm{chol}} L_{\mathrm{chol}}^T,
    \end{equation}
    caratterizzata da stabilit\`a numerica incondizionata e costo di appena $\frac{1}{3}n^3$ FLOPs.
\end{itemize}
Tale tematica sar\`a l'oggetto centrale del capitolo successivo.
